\documentclass{amsart}
% For dealing with references we use the comment environment
\usepackage{verbatim}
\newenvironment{reference}{\comment}{\endcomment}
%\newenvironment{reference}{}{}
\newenvironment{slogan}{\comment}{\endcomment}
\newenvironment{history}{\comment}{\endcomment}

% For commutative diagrams we use Xy-pic
\usepackage[all]{xy}

% We use 2cell for 2-commutative diagrams.
\xyoption{2cell}
\UseAllTwocells

% We use multicol for the list of chapters between chapters
\usepackage{multicol}

% This is generally recommended for better output
\usepackage{lmodern}
\usepackage[T1]{fontenc}

% For cross-file-references

% Package for hypertext links:
\usepackage{hyperref}

% For any local file, say "hello.tex" you want to link to please

% Theorem environments.
%
\theoremstyle{plain}
\newtheorem{theorem}[subsection]{Theorem}
\newtheorem{proposition}[subsection]{Proposition}
\newtheorem{lemma}[subsection]{Lemma}

\theoremstyle{definition}
\newtheorem{definition}[subsection]{Definition}
\newtheorem{example}[subsection]{Example}
\newtheorem{exercise}[subsection]{Exercise}
\newtheorem{situation}[subsection]{Situation}

\theoremstyle{remark}
\newtheorem{remark}[subsection]{Remark}
\newtheorem{remarks}[subsection]{Remarks}

\numberwithin{equation}{subsection}

% Macros
%
\def\lim{\mathop{\mathrm{lim}}\nolimits}
\def\colim{\mathop{\mathrm{colim}}\nolimits}
\def\Spec{\mathop{\mathrm{Spec}}}
\def\Hom{\mathop{\mathrm{Hom}}\nolimits}
\def\Ext{\mathop{\mathrm{Ext}}\nolimits}
\def\SheafHom{\mathop{\mathcal{H}\!\mathit{om}}\nolimits}
\def\SheafExt{\mathop{\mathcal{E}\!\mathit{xt}}\nolimits}
\def\Sch{\mathit{Sch}}
\def\Mor{\mathop{\mathrm{Mor}}\nolimits}
\def\Ob{\mathop{\mathrm{Ob}}\nolimits}
\def\Sh{\mathop{\mathit{Sh}}\nolimits}
\def\NL{\mathop{N\!L}\nolimits}
\def\CH{\mathop{\mathrm{CH}}\nolimits}
\def\proetale{{pro\text{-}\acute{e}tale}}
\def\etale{{\acute{e}tale}}
\def\QCoh{\mathit{QCoh}}
\def\Ker{\mathop{\mathrm{Ker}}}
\def\Im{\mathop{\mathrm{Im}}}
\def\Coker{\mathop{\mathrm{Coker}}}
\def\Coim{\mathop{\mathrm{Coim}}}

% Boxtimes
%
\DeclareMathSymbol{\boxtimes}{\mathbin}{AMSa}{"02}

%
% Macros for moduli stacks/spaces
%
\def\QCohstack{\mathcal{QC}\!\mathit{oh}}
\def\Cohstack{\mathcal{C}\!\mathit{oh}}
\def\Spacesstack{\mathcal{S}\!\mathit{paces}}
\def\Quotfunctor{\mathrm{Quot}}
\def\Hilbfunctor{\mathrm{Hilb}}
\def\Curvesstack{\mathcal{C}\!\mathit{urves}}
\def\Polarizedstack{\mathcal{P}\!\mathit{olarized}}
\def\Complexesstack{\mathcal{C}\!\mathit{omplexes}}
% \Pic is the operator that assigns to X its picard group, usage \Pic(X)
% \Picardstack_{X/B} denotes the Picard stack of X over B
% \Picardfunctor_{X/B} denotes the Picard functor of X over B
\def\Pic{\mathop{\mathrm{Pic}}\nolimits}
\def\Picardstack{\mathcal{P}\!\mathit{ic}}
\def\Picardfunctor{\mathrm{Pic}}
\def\Deformationcategory{\mathcal{D}\!\mathit{ef}}

\setlength{\emergencystretch}{3em}
\begin{document}
\title[Stacks Project, Tag 07PH]{346 --- A derivation detecting a non-pth-power in a finite-type domain}
\maketitle
\noindent Copyright (C) 2005--2025 Johan de Jong. Stacks Project authors. Source: \href{https://stacks.math.columbia.edu/tag/07PH}{Tag 07PH}.

\noindent Editorial difficulty note (not author text). Difficulty H2: The proof reduces by normalization and denominator scaling to a finite algebra over a power-series polynomial ring. A carefully chosen finite-index subfield preserves the nonzero differential, allowing a fraction-field derivation to be scaled back to the domain without killing the chosen element..

\noindent Editorial provenance note (not author text). History: excerpt from revision \texttt{a0444\allowbreak 6e57e\allowbreak c1fbc\allowbreak 252a8\allowbreak 71afc\allowbreak ec775\allowbreak 2fb28\allowbreak 07b14}, more-algebra.tex, lines 11859--11910. Reference presentation was adapted; original-author context and core support blocks were retained or added. The mathematical wording is unchanged. Licensed under GNU FDL 1.2 or later; no invariant sections or cover texts. See ../licenses/COPYING. Context blocks reproduce author definitions and supporting results; they are not additional counted problems.

\begin{definition}
\label{definition-p-basis}
Let $p$ be a prime number. Let $k \to K$ be an extension of fields
of characteristic $p$. Denote $kK^p$ the compositum of $k$ and $K^p$
in $K$.
\begin{enumerate}
\item A subset $\{x_i\} \subset K$ is called {\it p-independent
over $k$} if the elements $x^E = \prod x_i^{e_i}$ where
$0 \leq e_i < p$ are linearly independent over $kK^p$.
\item A subset $\{x_i\}$ of $K$ is called a
{\it p-basis of $K$ over $k$} if the elements
$x^E$ form a basis of $K$ over $kK^p$.
\end{enumerate}
\end{definition}

\begin{lemma}
\label{lemma-find-D}
Let $p$ be a prime number. Let $B$ be a domain with $p = 0$ in $B$.
Let $f \in B$ be an element which is not a $p$th power in the fraction
field of $B$. If $B$ is of finite type over a Noetherian complete
local ring, then there exists a derivation $D : B \to B$ such that $D(f)$
is not zero.
\end{lemma}

\begin{proof}
Let $R$ be a Noetherian complete local ring such that there exists
a finite type ring map $R \to B$. Of course we may replace $R$ by
its image in $B$, hence we may assume $R$ is a domain of characteristic
$p > 0$ (as well as Noetherian complete local). By Algebra, Lemma
\href{https://stacks.math.columbia.edu/tag/032D}{lemma complete local Noetherian domain finite over regular (Tag 032D)}
we can write $R$ as a finite extension of $k[[x_1, \ldots, x_n]]$ for some
field $k$ and integer $n$. Hence we may replace $R$ by $k[[x_1, \ldots, x_n]]$.
Next, we use
Algebra, Lemma \href{https://stacks.math.columbia.edu/tag/07NA}{lemma Noether normalization over a domain (Tag 07NA)}
to factor $R \to B$ as
$$
R \subset R[y_1, \ldots, y_d] \subset B' \subset B
$$
with $B'$ finite over $R[y_1, \ldots, y_d]$ and $B'_g \cong B_g$
for some nonzero $g \in R$. Note that $f' = g^{pN} f \in B'$ for some
large integer $N$. It is clear that $f'$ is not a $p$th power in
the fraction field of $B'$. If we can find a derivation
$D' : B' \to B'$ with $D'(f') \not = 0$, then
Lemma \href{https://stacks.math.columbia.edu/tag/07PE}{lemma derivation extends (Tag 07PE)}
guarantees that $D = g^MD'$ extends to $B$ for some $M > 0$. Then
$D(f) = g^MD'(f) = g^MD'(g^{-pN}f') = g^{M - pN}D'(f')$ is nonzero.
Thus it suffices to prove the lemma in case
$B$ is a finite extension of $A = k[[x_1, \ldots, x_n]][y_1, \ldots, y_m]$.

\medskip\noindent
Assume $B$ is a finite extension of
$A = k[[x_1, \ldots, x_n]][y_1, \ldots, y_m]$.
Denote $L$ the fraction field of $B$.
Note that $\text{d}f$ is not zero in $\Omega_{L/\mathbf{F}_p}$, see
Algebra, Lemma \href{https://stacks.math.columbia.edu/tag/031W}{lemma derivative zero pth power (Tag 031W)}.
We apply Lemma \href{https://stacks.math.columbia.edu/tag/07P5}{lemma power series ring subfields (Tag 07P5)} to find a subfield
$k' \subset k$ of finite index such that with
$A' = k'[[x_1^p, \ldots, x_n^p]][y_1^p, \ldots, y_m^p]$
the element $\text{d}f$ does not map to zero in $\Omega_{L/K'}$
where $K'$ is the fraction field of $A'$.
Thus we can choose a $K'$-derivation $D' : L \to L$
with $D'(f) \not = 0$. Since $A' \subset A$ and $A \subset B$ are
finite by construction we see that $A' \subset B$ is finite.
Choose $b_1, \ldots, b_t \in B$ which generate $B$ as an $A'$-module.
Then $D'(b_i) = f_i/g_i$ for some $f_i, g_i \in B$ with $g_i \not = 0$.
Setting $D = g_1 \ldots g_t D'$ we win.
\end{proof}
\end{document}
